\chapter{The Planar Dilatation Operator and Spin Chains}
\label{ch:volume-iii-planar-dilatation-operator-spin-chains}

\section{The SU(2) Scalar Sector}

In \(\mathcal N=4\) super-Yang--Mills, choose two complex adjoint scalars
\[
  Z=\frac1{\sqrt2}(\phi^1+\ii\phi^2),
  \qquad
  X=\frac1{\sqrt2}(\phi^3+\ii\phi^4).
\]
The \(SU(2)\) scalar sector is spanned, at the classical level, by
single-trace words
\begin{equation}
  \operatorname{Tr}(A_1A_2\cdots A_L),
  \qquad
  A_\ell\in\{Z,X\}.
  \label{eq:su2-sector-single-trace-word}
\end{equation}
The integer \(L\) is the length of the word.  The operator
\[
  \operatorname{Tr}Z^L
\]
is a half-BPS primary with \(\Delta=L\).  Replacing some of the \(Z\)'s by
\(X\)'s creates excitations above this protected reference operator.

At large \(N\), the one-loop dilatation operator preserves the \(SU(2)\)
sector.  It acts on neighboring sites of the chain.  Let
\[
  P_{\ell,\ell+1}
\]
be the operator that exchanges the fields at sites \(\ell\) and \(\ell+1\),
with \(\ell+1\) understood modulo \(L\).  The planar one-loop anomalous
dimension operator is
\begin{equation}
  \gamma^{(1)}
  =
  \frac{\lambda}{8\pi^2}
  \sum_{\ell=1}^{L}
  \left(1-P_{\ell,\ell+1}\right).
  \label{eq:su2-one-loop-dilatation}
\end{equation}
Thus the full scaling dimension in this sector is
\[
  \Delta
  =
  L+\gamma^{(1)}+O(\lambda^2).
\]
Equation \eqref{eq:su2-one-loop-dilatation} is the Hamiltonian of the
ferromagnetic Heisenberg \(XXX_{1/2}\) spin chain.

\begin{figure}[H]
  \centering
  \resizebox{0.90\textwidth}{!}{%
  \begin{tikzpicture}[x=1cm,y=1cm,every node/.style={font=\scriptsize}]
    \tikzset{arrow/.style={-{Latex[length=2mm]},line width=0.85pt}}
    \foreach \x/\lab in {0/{Z},0.9/{Z},1.8/{X},2.7/{Z},3.6/{X},4.5/{Z},5.4/{Z}} {
      \draw[thick,fill=blue!5] (\x-0.28,-0.28) rectangle (\x+0.28,0.28);
      \node at (\x,0) {\(\lab\)};
    }
    \draw[orange!85!black,line width=0.95pt] (1.50,-0.48) rectangle (3.00,0.48);
    \node[orange!85!black] at (2.25,0.82) {\(1-P_{\ell,\ell+1}\)};
    \draw[arrow] (6.05,0)--(7.25,0);
    \node[align=center] at (8.20,0.35)
      {nearest-neighbor\\spin-chain Hamiltonian};
    \node[align=center,blue!65!black] at (2.70,-1.10)
      {\(Z\) is spin up, \(X\) is spin down};
  \end{tikzpicture}
  }
  \caption{In the planar \(SU(2)\) sector, the one-loop anomalous dimension
  operator is a nearest-neighbor spin-chain Hamiltonian.}
  \label{fig:su2-sector-nearest-neighbor-chain}
\end{figure}

\section{Magnon Eigenstates}

On a long chain, an \(X\) inserted among \(Z\)'s behaves as a magnon.  Let
\(\ket{\ell}\) denote the state with the impurity \(X\) at site \(\ell\):
\[
  \ket{\ell}
  =
  \operatorname{Tr}(Z^{\ell-1}XZ^{L-\ell}).
\]
A momentum eigenstate is
\begin{equation}
  \ket p
  =
  \sum_{\ell=1}^{L}
  \ee^{\ii p\ell}\ket{\ell},
  \qquad
  p=\frac{2\pi n}{L},
  \quad n\in\mathbb Z.
  \label{eq:single-magnon-momentum-state}
\end{equation}
Acting with \eqref{eq:su2-one-loop-dilatation} gives
\begin{equation}
  \gamma^{(1)}(p)
  =
  \frac{\lambda}{2\pi^2}
  \sin^2\frac p2.
  \label{eq:one-loop-magnon-dispersion}
\end{equation}
Cyclicity of the trace imposes the level-matching condition that the total
spin-chain momentum of a physical single-trace state is zero modulo
\(2\pi\).  A nonzero-momentum single magnon is therefore best regarded as the
elementary excitation used to build multi-magnon states; physical closed-chain
states have total momentum zero.

\section{Protected Vacuum and Missing Diagrams}

The operator \(\operatorname{Tr}Z^L\) is annihilated by every term
\[
  1-P_{\ell,\ell+1}
\]
because adjacent \(Z\)'s are identical.  Hence
\[
  \gamma^{(1)}\operatorname{Tr}Z^L=0.
\]
This implements the half-BPS protection of \(\operatorname{Tr}Z^L\) at one
loop.  In a Feynman-diagram computation, self-energy and exchange diagrams
separately contribute to the anomalous dimension matrix; the sum is fixed so
that the protected vacuum has zero anomalous dimension.  The spin-chain form
\eqref{eq:su2-one-loop-dilatation} packages that cancellation into the local
operator \(1-P\).

\begin{figure}[H]
  \centering
  \resizebox{0.92\textwidth}{!}{%
  \begin{tikzpicture}[x=1cm,y=1cm,every node/.style={font=\scriptsize}]
    \tikzset{arrow/.style={-{Latex[length=2mm]},line width=0.85pt}}
    \draw[thick] (0,0) ellipse (1.45 and 0.55);
    \foreach \ang/\lab in {20/{Z},70/{Z},120/{X},175/{Z},230/{Z},290/{X},335/{Z}} {
      \node at ({1.45*cos(\ang)},{0.55*sin(\ang)}) {\(\lab\)};
    }
    \draw[magenta,line width=0.95pt] ({1.45*cos(120)},{0.55*sin(120)}) circle (0.18);
    \draw[magenta,line width=0.95pt] ({1.45*cos(290)},{0.55*sin(290)}) circle (0.18);
    \node[magenta] at (0,1.25) {magnons};

    \draw[arrow] (2.0,0)--(3.1,0);
    \node[align=center] at (2.55,0.50) {diagonalize\\\(\gamma\)};

    \begin{scope}[shift={(5.15,0)}]
      \draw[->,thick] (-1.25,-0.65)--(1.35,-0.65) node[right] {\(p\)};
      \draw[->,thick] (-1.15,-0.78)--(-1.15,1.05) node[above] {\(\gamma\)};
      \draw[blue!65!black,line width=1pt,domain=-1.0:1.0,samples=80]
        plot (\x,{0.75*sin(90*\x)*sin(90*\x)-0.55});
      \node[blue!65!black] at (0.15,0.75)
        {\(\frac{\lambda}{2\pi^2}\sin^2\frac p2\)};
    \end{scope}
  \end{tikzpicture}
  }
  \caption{Magnons are impurities propagating on the protected trace
  \(\operatorname{Tr}Z^L\).  Their one-loop energy is the eigenvalue of the
  planar dilatation operator.}
  \label{fig:magnon-dispersion-chain}
\end{figure}

\section{All-Loop Dispersion and Factorization}

The one-loop magnon dispersion is the first term in an expression constrained
by the centrally extended \(SU(2|2)\) symmetry of the planar theory.  In the
normalization used above, the all-loop single-magnon dispersion relation is
\begin{equation}
  \Delta-J
  =
  \sqrt{
    1+\frac{\lambda}{\pi^2}\sin^2\frac p2
  },
  \label{eq:all-loop-magnon-dispersion}
\end{equation}
where \(J\) is the \(U(1)\subset SU(4)_R\) charge carried by the \(Z\)'s.
Expanding \eqref{eq:all-loop-magnon-dispersion} at small \(\lambda\) gives
\[
  \Delta-J
  =
  1+\frac{\lambda}{2\pi^2}\sin^2\frac p2
  +O(\lambda^2),
\]
which agrees with \eqref{eq:one-loop-magnon-dispersion} after subtracting the
bare impurity dimension.

Planar integrability asserts that multi-magnon scattering factorizes into
two-body scattering.  The two-body S-matrix acts on the tensor product of
magnon internal spaces and is constrained by symmetry, unitarity, crossing,
and analyticity.  The resulting Bethe equations determine the planar
single-trace spectrum in sectors where wrapping effects are absent or have
been incorporated by finite-size corrections.

The conceptual role of the spin chain is therefore precise: it is the
diagonalization problem for the dilatation operator acting on local
single-trace operators in the planar conformal theory.  Its Hilbert space,
Hamiltonian, and scattering data are the large-\(N\) representation of the
operator-mixing problem of the conformal field theory.
